%!TEX program = uplatex
\RequirePackage{plautopatch}
\documentclass[uplatex,dvipdfmx]{jlreq}
\usepackage{amsmath,amssymb}
\usepackage{enumerate}
\pagestyle{empty}

\begin{document}

%======================================================================
% 講演１：泉屋周一「一階偏微分方程式への応用」
%======================================================================
\begin{center}
{\large\bfseries 一階偏微分方程式への応用}

\hspace*{\fill}{\large\bfseries 泉屋 周一（北大・理）}
\end{center}

\medskip

線形の偏微分方程式に対しては、弱解として超関数の理論が大変良く機能することが、
今世紀前半の微分方程式論における主要な研究成果の一つであった。
しかし、非線形偏微分方程式の場合、現在までのところ統一的な取り扱いは存在しない
（ように思われる）。
たとえ一階の偏微分方程式においてでさえも、準線形の場合と非線形の場合には取り扱いが異なる。
一方、一階の場合には、幾何学的方法が前世紀の偉大な数学者達
（リー、ダルブー、グルサー等）によって、ほぼ完全な形に研究されてきた。
しかし、幾何学的理論と応用解析学的研究には明確なギャップが存在する。
幾何学的理論におけるひとつの優位性は解の多価性を許容することにより
微分方程式を部分多様体論として記述し、解の存在や一意性の導出を容易にすることである。
しかし、現実への応用を重視する応用解析学においては、解が一価であることを要請する場合が多い。
そのために、解の微分可能性という重要な性質も捨て去るのである。

ここでは、一階の偏微分方程式のなかで、ハミルトン・ヤコビ方程式について、
特性曲線の方法（幾何学的理論）で解かれて得られる多価解（ルジャンドル解）と
1983年ごろにクランダールとリオンスによって導入された粘性解との関係について述べる。
ハミルトン・ヤコビ方程式に対して初期条件が滑らかなコーシー問題を特性曲線の方法で解くと、
十分短い時間のうちはその独立変数に関するヤコビ行列が退化しないので
滑らかな逆関数を求めることができ、それにより滑らかな一価解（古典解）が
構成できることは偏微分方程式の入門書には良く書かれている事実である。
しかし、そのヤコビ行列が退化した場合については、ほとんど記述が無い。
それ以降は、全く別の方法により弱解の存在や一意性を示しているのが応用解析的立場であり、
また一価性を捨て去り、特性曲線のばせるだけのばして解を部分多様体として捉えるのが幾何学的立場である。
さて、ヤコビ行列が退化する点はいわゆる可微分写像としての特異点である。
したがって、そのような点の回りでの解の振る舞いを研究することは、特異点論の範疇に属する。

ここでは、特性曲線の方法で解かれた多価解がルジャンドル開折というある種の
ルジャンドル部分多様体のクラスに属すること、
そしてその特異点の分類が与えられることを解説する。
その後、そこからどのようにして一価な粘性解を構成するかについて解説する。
特に、ハミルトン関数が凸（または凹）の場合は多価解のある連続な分枝を選んで構成できるが、
ハミルトン関数が変曲点を持つ場合は新たに特性曲線を構成しなおす必要があることを解説する。
また、この凸（または凹）の場合の理論の応用として、
多価解を逆に粘性解から数値解析的に構成する試みについて述べる。
この方法は、従来の特性曲線の方法を直接適用するよりずっと安価な計算で可能なことが期待されている。

\newpage

%======================================================================
% 講演２：泉屋周一「微分幾何学への応用」
%======================================================================
\begin{center}
{\large\bfseries 微分幾何学への応用}

\hspace*{\fill}{\large\bfseries 泉屋 周一（北大・理）}
\end{center}

\medskip

特異点論を微分幾何学へ応用する試みは1970年代の Thom による Cut locus の研究に遡る。
しかし、変分問題などを考えればそれは数学上の自然な流れであり、
Thom が特異点論という立場から再定式化したものである。
この事に関連して、それ以前よりモース理論等のいちじるしい成果があるが、
Thom の考えは単独のジェネリックな関数の特異点を考える（モース理論などはそうです）のではなく、
関数族に付随した判別集合や分岐集合の特異点を通して幾何学的不変量を研究することにあり、
それはとりもなおさず「初等カタストロフィー」の考えである。

実際、Thom はユークリッド空間内の部分多様体上の距離２乗関数や高さ関数を考えることにより、
その部分多様体の古典的不変量が研究できることを示唆した。
この考えを具体的に実行したのは Porteous であり、その後の発展は彼の著書の中に見られる。
これらの関数を考えることは微分幾何学的にも自然なことであり、
モース関数も具体的にはこれらの関数で与えられているし、
これらの関数を使うことにより tight（taut）はめ込みの研究などが行われている。

一方、Thom の考えはユークリッド空間内の部分多様体以外にも適用可能なことは容易に想像がつく。
しかし、各種の幾何学に対するこの立場からの研究は近年まで全く適用された例は見あたらない。
ここでは、アファイン微分幾何学やローレンツ微分幾何学、双曲幾何学などへの特異点論を応用することにより、
得られる様々な幾何学的不変量とその幾何学的意味について解説する。
また、偏微分方程式の解として現れる曲面の特異点と微分幾何学的性質の関係に関する
研究の紹介も行いたい。

\newpage

%======================================================================
% 講演３：石川剛郎「『応用特異点論』の基礎」
%======================================================================
\begin{center}
{\large\bfseries 「応用特異点論」の基礎}

\hspace*{\fill}{\large\bfseries 石川 剛郎（北大・理）}
\end{center}

\medskip

いろいろな分野で、古典的にあいまいに扱われてきたこと、
難しすぎてとても手に負えなかった問題などが、
写像の特異点を調べる組織的方法によって、いま解かれつつあります。
横断性、安定性、確定性、普遍性といった、特異点論において標準化された処方箋が、
われわれを解決へと導いてくれるのです。
哲学不在の現代、いかに現象を認識するか、
その根源的視座のひとつとして、特異点論は発展を続けています。

写像の特異点論とは何か。特異点論の基本的な考え方はどのようなものか。それは何をめざしているのか。
このような問いに、象徴的な具体例によって可能な限り答えたいと思います。

私（石川）の講演では、理論のプロトタイプとして、
平面から空間への写像や、平面から平面への写像に現われる特異点がどのようなものか説明します。
また、モース関数や、もっと退化した関数、
あるいは超曲面の特異点を、どのように扱い分類するかという話をします。
数学的ディテールにあまりこだわらず、この素材をもとに、
動機づけや理念を中心に、わかりやすく話すよう努めます。

後半では、特異点論の言葉づかいや、シンプレクティック幾何や接触幾何などに触れつつ、
ラグランジュ・ルジャンドル特異点論の説明をします。
ラグランジュ・ルジャンドル特異点論は、進化した数学的カタストロフ理論というべきものであり、
様々な応用が期待されていると同時に、理論自身、いまも日進月歩発展を続けています。

私の講演は、「応用特異点論」でいうと、
第３章から第７章までの基礎的な部分の概説といえます。

では、中央大学でお会いしましょう。

\newpage

%======================================================================
% 講演４：佐伯修「微分位相幾何学への応用」
%======================================================================
\begin{center}
{\large\bfseries 微分位相幾何学への応用}

\hspace*{\fill}{\large\bfseries 佐伯 修（広島大・理）}
\end{center}

\medskip

本講演では、微分可能写像の特異点論の微分位相幾何学への応用（のほんの一部）について、
最近の結果を中心にいくつか述べたいと思う。

写像の特異点論の応用は、微分位相幾何学の中に数多く見ることができる。
たとえば、今世紀前半にホイットニーは、どんな $n$ 次元多様体も、
$2n-1$ 次元ユークリッド空間にはめ込んだり、$2n$ 次元ユークリッド空間に埋め込んだりできる
という定理を示したが、これは（微分位相幾何学的な意味での）特異点の解消定理と理解することができる。
また、変分問題において重要な働きをする多様体上のモース理論は、
関数の非退化臨界点に関するモースの補題という特異点論の定理を
多様体の位相的構造を調べるのに使ったと解釈することができる。
さらに、最近はやりのバシリエフ不変量の理論も、写像空間の中で、
特異点を持つ写像の空間の補空間のトポロジーを調べるという意味で、
特異点理論の応用の一つであると言える。

このように話すべき話題は多いわけであるが、
本講演では、ある限られた簡単な特異点しか持たない微分可能写像に着目し、
与えられた多様体に対して、そうした写像を許容する・しないという事実が、
その多様体の位相構造、ひいては微分構造と深くかかわるという
最近明らかにされた事実について解説してゆきたいと思う。
たとえば、ある種の写像を許容することと、大域変分学において重要な役割を果たした
ルシュターニーク・シュニーレルマンカテゴリーが密接に関連すること、
解析的手法を用いて定義される４次元多様体のサイバーグ・ウィッテン不変量を用いて、
複素解析曲面上のある種の写像の非存在が示せること、などについて言及したい。

\begin{center}\textsc{参考文献}\end{center}

\begin{enumerate}[{[1]}]
  \setlength{\itemsep}{3pt}
  \item[{[1]}] 佐伯修，微分可能写像の大域的特異点理論の現状と展望，
        数学 \textbf{48} (1996), 385--399.
  \item[{[2]}] 佐伯修，微分位相幾何と特異点，数理科学，1998年6月号，43--49.
\end{enumerate}

\end{document}
